%%%                          Súbor ZFSme.tex                             %%%
% Obsahuje vlastné makrá pre spracovanie  kníh  z  rtf  súborov  zo  stránok
% http://zlatyfond.sme.sk, ktoré boli  odladené  pre  knihu  Hviezdoveda  od
% Gustáva Reussa. Pre spracovanie iných  kníh  bude  pravdepodobne  potrebné
% doplniť niektoré makrá. Pre vlastné makrá je okrem samotného systému plain
% pdfTeX potrebný aj CSTeX a súbor makier OPmac.

%%%            Príklad konfigurácie v sparacovávanom súbore .tex          %%%
%% V súbore ZfSME.tex nechať všetky možnosti zakomentované, nastavovať     %%
%% konfiguráciu až v súbore, ktorý budeme spracovávať TeX-om               %%

%% Úvodné nastavenie CSTeX a OPmac kde by malo byť aspoň toto:
%\shyph % zapína slovenské delenie slov vo formáte csplain/pdfcsplain
%\input opmac % používa sa verzia OPmac Apr. 2020 (alebo novšia)
%\raggedbottom           % nezarovnaný spodný okraj strany

%% Potlačenie výskytu sluru; odkomentovať až po odladení výstupu
%\overfullrule=0pt

%% Nastavenie klikacích odkazov (funkčné len pri výstupe do pdf),
%% vo formáte {interný}{externý -- napr. www adresa} odkaz. 
%% Po zakomentovaní oboch riadkov nebudú žiadne klikacie odkazy funkčné.
%\hyperlinks{\Blue}{\Green}    % nastavenie farebných klikacích odkazov
%\hyperlinks{}{}               % nastavenie č/b klikacích odkazov

%% \Capinsert První písmeno bude větší, tj. P v tomto případě...
%% https://petr.olsak.net/opmac-tricks.html#capinsert
%% Potrebné odladiť na: 
%% - viac odstavcov zasahujúcich prvé písmeno: \nli, \nlii
%% - automaticky posunúť nahor prvé písmeno s diakritikou: \VshiftFont
%% - pri použití závesnej interpunkcie na začiatku riadku zasahuje do písmena
%% - zmena fontu vyžaduje zmenu parametrov, hlavne s diakritikou: \CapIskip,
%%   \CapIIskip pridáva medzeru nad písmeno s \CapInsert, CapIInsert
%% po odkomentovaní zapnutie funkcie \Capinsert
%\def\IfCapInsert{1} % 1 - len pre predhovory, 2 - pre všetky kapitoly
%\def\CapNumLines{2}			  % počet riadkov písmena (2 alebo 3)
%\newdimen\ptex       \ptex=.1ex          % jednotka závislá na ex 
%\newdimen\VshiftFont \VshiftFont=0\ptex  % vertikálny posun pre písmeno
%\def\Capprefix{\localcolor\Red}          % makro provedené před písmenem 
%\def\Capprefix{\localcolor\Grey}         % makro provedené před písmenem 
%
%% Zmena použitých základných fontov, namiesto prednastaveného:
%\def\FontType{Computer Modern vo variante \CS-font} \let\CMCS=\FontType
%% Pôvodné fonty z inšalácie \CS\TeX (majú spoločný rez pre \it a \sl),
%% (odkomentovať len jeden z nasledovných 4 riadkov):
%\input ctimes \def\FontType{Times Roman} 
%\input cbookman \def\FontType{Bookman} \VshiftFont=2\ptex
%\input cncent \def\FontType{New Century} \VshiftFont=4.4\ptex
%\input cpalatin \def\FontType{Palatino} \VshiftFont=4\ptex 
%% Nové fonty z inšalácie \CS\TeX (rez pre \sl je použitý z CS-CM,
%% pokiaľ nie je aktívna RS \let\sl=\it na konci riadku);
%% (odkomentovať len jeden z nasledovných 4 riadkov):
%\input lmfonts \def\FontType{Latin Modern} % odporúčaná náhrada za CM
%\input cs-charter \def\FontType{Charter}  \let\sl=\it \VshiftFont=3\ptex
%\input cs-antt \def\FontType{Antykwa To\-ru\'nska}  \let\sl=\it % *
%\input cs-polta \def\FontType{Antykwa Pó\l taw\-skiego}  \let\sl\it % *
%% * v niektorých systémoch je potrebné doinštalovať tieto fonty do TeXLive
%% (napr. balík texlive-fonts-extra pre Debian/Ubuntu)

%% Nastavenie formátu obrázkov (png farebný/png šedý/eps)
%\def\ObrPdf{0}
%\def\ObrPdf{1} % zakomentovať riadok pre použitie farebných png obrázkov
%%\def\ObrPdf{2} \input epsf % odkomentovať riadok pre výstup do dvi (už nie
                             % je potrebné, príslušné makrá si to nastavia)

%% Nastavenie grafických súborov pre logo Zlatý fond SME
%\def\LogoZfSMEbc{zlaty_fond.png} % veľké farebné logo ZfSME pre PDF výstup
%\def\LogoZfSMEbg{zlaty_fond-g.png} % veľké č/b logo ZfSME pre PDF výstup
%\def\LogoZfSMEbe{zlaty_fond.eps} % veľké farebné logo ZfSME pre DVI výstup
%\def\LogoZfSMEsc{zlaty_fond_maly.png} % malé farebné logo ZfSME päty strany
%\def\LogoZfSMEsg{zlaty_fond_maly-g.png} % malé č/b logo ZfSME päty strany

%% Nastavenie zobrazenia loga Zlatý fond SME v strede päty strany:
%\def\LogoZF{-1} % namiesto loga čísla strán v strede päty (pre e-book)
%\def\LogoZF{0} % bez zobrazenia loga
%\def\LogoZF{1} % zobarzenie textu loga
%\def\LogoZF{2} % zobrazenie grafického loga (len pre výstup do pdf)
%\def\ZfSME{{\rm Zlatý fond SME}} % definovanie textu a formátu loga

%% Nastavenie tlače záhlavia strany:
%\def\ZahlavieTlac{0} 
%\def\ZahlavieTlac{1} % netlačí sa po zakomentovaní riadku

%% Potlačenie tlače záhlavia a päty strany na prázdnych stranách:
%\def\EmptyHLFL{0} % tlačí sa po odkomentovaní riadku (neodporúča sa)

%% Nastavenie hlavičky pre motto:
%\def\MottoH{Motto:} % zobrazí sa v hlavičke motta po odkomentovaní

%% Zapnutie automatického číslovania kapitol a pod/sekcií po odkomentovaní
%\def\NoNum{0} % nastaví sa po odkomentovaní riadku

%% Zapnutie linkov na strany pre index po odkomentovaní
%\def\IndexLink{1}
%% Zapnutie delenia podľa začiatočného písmena pre index po odkomentovaní
%\def\IndexChar{1} 

%% Nastavenie formátu zlomku:
%\def\FRformat{0} % formát zlomkovej čiary  /
%\def\FRformat{1} % formát zlomkovej čiary ---

%% Nastavenie minimálnej veľkosti šírky strany, pri ktorej sa bude text
%% zarovnávať do bloku bez zväčšenia medzier:
%\def\MinHSize{100mm}
%\emergencystretch=2em % nastavenie väčších medzier bez podmienky \MinHSize
%% Nastavenie minimálnej veľkosti šírky strany, pri ktorej sa bude text
%% zarovnávať do bloku (pri užšej strane bude text zarovnaný vľavo):
%\def\MinHsize{80mm}

%% Nastavenie zarovnania vľavo pre sekciu a popodsekciu:
%\def\SecLeft{0}
%\def\SecLeft{1} % zarovnanie pod/pod/sekcie do bloku po zakomentovaní riadku 

%% Nastavenie min. veľkosti ľavého + pravého okraja pre závesnú interpunkciu:
%\def\MinOkraj{9cm} % odkomentovať pre vypnutie aj pre knižné formáty

%% Doplnenie veľkosti strán/formátov pre A6 a e-booky:
%\sdef{pgs:a6}{(105,148)mm} 
%\sdef{pgs:eb}{(96,156)mm}  \sdef{pgs:ebr}{(96,96)mm}
%\sdef{pgs:ebsr}{(78,78)mm} % 78mm je najmenší rozmer pre tabuľku Miery a váhy
%\sdef{pgs:názov}{(<šírka>,<výška>)} % vlastný formát
%% Nastavenie základnej veľkosti použitých fontov a veľkosti strany,  po zako-
%% mentovaní všetkých riadkov v tejto časti sa použije nastavenie pre csplain:
%% strana veľkosti A4, okraje 2,54 cm, veľkosť fontu 10 pt, riadkovanie 12 pt.
%% Font veľkosti 11 pt, riadkovanie 12,5 pt, vhodné pre väčšiu veľkosť strany:
%\typosize[11/12.5] \VshiftFont=1.25\VshiftFont
%\margins/2 a4 (20,37,27,52)mm % formát A4, P/Ľ zrkadlo, pre a=1/4 šírky strany
%\margins/2 a4 (17,30,24,42)mm % formát A4, P/Ľ zrkadlo, pre a=1/5 šírky strany
%\margins/1 b5 (24,24,24,24)mm % nastavenie formátu B5 a všetky okraje na 24mm
%\margins/2 b5 (14,26,20,35)mm % formát B5, P/Ľ zrkadlo, pre a=1/5 šírky strany
%\margins/1 a5 (18,18,18,18)mm % nastavenie formátu A5 a všetky okraje na 18mm
%\margins/2 a5 (14,22,18,28)mm % formát A5, P/Ľ zrkadlo, pre a=1/5 šírky strany
%% Nastavenia strany pre e-book (zariadenia formátu A6, alebo vlastný formát)
%\margins/1 a6 (3,3,3,11)mm % nastavenie formátu A6, okraje na 3mm, spodný 11mm
                            % (zobrazí sa aj číslo strany v strede päty)
%\margins/1 eb (3,3,3,11)mm % nastavenie formátu na ebook, okraje na 3 a 11mm 
                            % (zobrazí sa aj číslo strany v strede päty)
%\margins/1 ebsr (1,1,1,1)mm %nastavenie formátu na ebook, všetky okraje na 1mm
                            %(v tomto prípade sa nebudú zobrazovať čísla strán)

%% Posun titulu na titulnej strane na cca zlatý stred
%\def\TitSkip{.05} % prednastavená hodnota je 0,1 výšky strany
%\ifdim\vsize < \MinHsize % len pre ebsr, už musí byť známa veľkosť \vsize
%\def\TitSkip{-.1} % odkomentovať len pre viac mien, alebo dlhý názov
%\fi

%% Načítanie vlastných makier pre Zlatý fond Sme:
%\ifx\ZFroot\undefined \def\ZFroot{} \fi % ak nie je iná cesta k ZfSME.tex
%\input \ZFroot ZfSME
%% Upresnenie umiestnenia veľkých začiatočných písmen (až po načítaní ZfSME)
%\declCap J {1:26;2,1} %odkomentovať pre Palatino a AntykwaTorunska (2 riadky)
%\declCap J {1:38;2,0,0} %odkomentovať pre Palatino a AntykwaTorunska (3 riad.)

%% Nastavenie informácií o titule a autorovi; v prípade dlhého názvu alebo mena
%% je možné použiť \nl na oddelenie riadkov v parametri \Titul a \Autor napr.
%% \def\Autor{Meno \nl Priezvisko}
%\def\Titul{Titul}
%\def\Autor{Meno Priezvisko}
%% Nastavenie informácií v PDF - viď TPP, strany 94 až 95:
%\input pdfuni % umožní vkladať PDF informácie v kódovaní UTF-16BE Unicode
%\pdfunidef\tmp{\Autor}    \pdfinfo{/Author (\tmp)}
%\pdfunidef\tmp{\Titul}    \pdfinfo{/Title (\tmp)}

%%% Vlastné makrá (začínajú sa obyčajne veľkým písmenom) %%%

% V prípade že sa nepoužíva TeXonWeb -- https://tex.mendelu.cz pomocou
%\input TeXonWeb.tex % je potrebné definovať makro \ToWp :
\ifx\ToWp\undefined \def\ToWp{} \fi

% Pre staršie verzie CSTeX nemusí byť \elq definované
\ifx\elq\undefined \chardef \elq 96 \fi

% Prednastavené hodnoty, pokiaľ neboli definované v konfigurácii:
\ifx\ZFroot\undefined \def\ZFroot{} \fi % ak nie je iná cesta k ZfSME.tex
\ifx\ZFbin\undefined \def\ZFbin{} \fi % ak nie je iná cesta k zfsme-rtf2tex
\ifx\FontType\undefined % definícia pokračuje na nasledujúcom riadku:
\def\FontType{Computer Modern vo variante \CS-font} \let\CMCS=\FontType \fi
\ifx\IfCapInsert\undefined \def\IfCapInsert{0} \fi
\ifx\CapNumLines\undefined \def\CapNumLines{2} \fi
\ifx\Capprefix\undefined \def\Capprefix{}      \fi
\ifx\ObrPdf\undefined \def\ObrPdf{0} \fi
\ifx\LogoZF\undefined \def\LogoZF{0} \fi
\ifx\LogoZfSMEbc\undefined \def\LogoZfSMEbc{zlaty_fond.png} \fi
\ifx\LogoZfSMEbg\undefined \def\LogoZfSMEbg{zlaty_fond-g.png} \fi
\ifx\LogoZfSMEbe\undefined \def\LogoZfSMEbe{zlaty_fond.eps} \fi
\ifx\LogoZfSMEsc\undefined \def\LogoZfSMEsc{zlaty_fond_maly.png} \fi
\ifx\LogoZfSMEsg\undefined \def\LogoZfSMEsg{zlaty_fond_maly-g.png} \fi
\ifx\ZfSME\undefined \def\ZfSME{{\it Zlatý fond SME}} \fi
\ifx\ZahlavieTlac\undefined \def\ZahlavieTlac{0} \fi
\ifx\EmptyHLFL\undefined \def\EmptyHLFL{1} \fi
\ifx\FRformat\undefined \def\FRformat{0} \fi
\ifx\MinHSize\undefined \def\MinHSize{100mm} \fi
\ifx\MinHsize\undefined \def\MinHsize{80mm} \fi
\ifx\SecLeft\undefined \def\SecLeft{0} \fi
\ifx\MinOkraj\undefined \def\MinOkraj{9mm} \fi
\ifx\TitSkip\undefined \def\TitSkip{.1} \fi
\ifx\NoNum\undefined \def\NoNum{1} \fi % potlačenie číslovania kapitol a sekcií
\ifx\MottoH\undefined \def\MottoH{} \fi
\ifx\IndexChar\undefined \def\IndexChar{0} \fi
\ifx\IndexLink\undefined \def\IndexLink{0} \fi
\ifx\Titul\undefined \def\Titul{Titul} \fi
\ifx\Autor\undefined \def\Autor{Autor} \fi
\ifx\ChapRow\undefinned\def\ChapRow{0}\fi
\ifx\SecRow\undefinned\def\SecRow{0}\fi
\ifx\SeccRow\undefinned\def\SeccRow{0}\fi
\ifx\KindleHR\undefined \def\KindleHR{} \fi
\isdefined{ptex }\iftrue\else\newdimen\ptex\fi % ak rozmer \ptex nie je
                                               % definovaný v konfigurácii 
% Definovanie makier pre logo METAFONT a XeTeX
\font\logo=logo10 %scaled\magstephalf
\font\logosl=logosl10 %scaled\magstephalf
\regfont\logo % Umožňuje automatické zväčšovanie písma logo v texte
% font logo nemá znak -, preto sa používa \discretionary{\rm -}{}{}:
\def\mf{{\logo META\discretionary{\rm -}{}{}FONT}} % logo METAFONT
\def\mfsl{{\logosl META\discretionary{\rm -}{}{}FONT}} % sklonená verzia 
% \reflectbox pre pdfTeX -- viď použitie OPmac makra \pdfscale, TPP str. 100 
\def\reflectbox#1{\phantom{#1}\pdfsave \pdfscale{-1}{1}\rlap{#1}\pdfrestore}
% veľkosti kernov pre logo \XeTeX viď 
% https://tex.stackexchange.com/questions/558031/xetex-logo-with-reversed-e
\def\XeTeX{X\lower.5ex\hbox{\kern-.125em\reflectbox{E}}\kern-.1667em\TeX}

% Odstránkovanie bez podmienok:
\def\NewPage{\ifdim\pagetotal < .5\vsize \null\fi \vfill\eject} % TBN str. 240
% Odstrákovanie len v prípade, že číslo strany je: 
% párne:
\def\NewPageIfEven{{\ifodd\pageno \else \ifnum\EmptyHLFL=1
                    \footline={\hfil} \headline={\hfil} \fi
                    \NewPage \fi }}
% nepárne:
\def\NewPageIfOdd{{\ifodd\pageno \ifnum\EmptyHLFL=1
                   \footline={\hfil} \headline={\hfil} \fi
                   \NewPage \fi }}

% Kontrola nevhodnej kombinácie výstupu do PDF a \ObrPdf
\ifpdftex \ifnum\ObrPdf=1 \else \def\ObrPdf{0} \fi
\else \ifnum\ObrPdf=2 \else \def\ObrPdf{2} \input epsf \fi\fi
% Kontrola nevhodnej kombinácie \LogoZF a \ObrPdf
\ifnum\LogoZF=2 \ifnum\ObrPdf=2 \def\LogoZF{1} \fi\fi
% Kontrola prítomnosti súborov s malým logom ZfSME
% (ak nie sú, použije sa textové logo ZfSME)
\def\InsPicFooter#1{\openin\testin=\ZFroot #1
\ifeof\testin\opwarning{The ‘#1’ file doesn’t exist..}%
\ifnum\LogoZF=2 \def\LogoZF{1} \fi \else
\closein\testin\pdfximage height24pt {\ZFroot #1}\fi}
\ifnum\ObrPdf = 0 \InsPicFooter{\LogoZfSMEsc}\fi % načítanie farebného loga
\ifnum\ObrPdf = 1 \InsPicFooter{\LogoZfSMEsg}\fi % načítanie č/b loga

% Nastavenie stránkovania a pdf záložiek pre e-booky
\ifnum\LogoZF = -1 % pre e-booky sa negenerujú prázdne strany
\def\NewPage{\vfill\eject} % bez \null, po \endmulti môže vytvoriť stranu 
\def\NewPageIfEven{} \def\NewPageIfOdd{}
\def\Outlines{\insertoutline{Obsah} \outlines{0}} % záložky v pdf pre e-booky
\else \def\Outlines{} \fi
% Nastavenie formátu päty bežnej strany:
\def\PageNumbersRoman{% vhodné pre obsah a úvodné strany
\footline={\hss\tenrm\folio\hss}}
\ifnum \LogoZF = -1
% verzia bez loga, s umiesnením č. strany v strede päty (vhodné pre e-book)
\let\PageNumbers=\PageNumbersRoman
% nasledovné rozloženia sú vhodné pre obojstrannú tlač:
\else \ifnum \LogoZF = 0
% verzia bez loga, s umiestnením č. strany (vľavo pre párnu a vpravo 
% pre nepárnu stranu (hfill je silnejšia "pružina" ako hfil):
\def\PageNumbers{%
\footline={\tenrm \ifodd\pageno \hfill \fi \folio \hfil}}
\else \ifnum \LogoZF = 1
% verzia s nápisom Zlatý fond SME v strede päty strany
\def\PageNumbers{%
\footline={\tenrm \ifodd\pageno \hphantom{\folio}\hfil\ZfSME\unskip\hfil \folio
                        \else
                        \folio \hfil \ZfSME\unskip\hfil\hphantom{\folio}\fi}}
\else \ifnum \LogoZF = 2
% verzia s obrázkom Zlatý fond SME v strede päty strany
\mathchardef\ZfSMEi=\pdflastximage % TPP, str. 101
\def\PageNumbers{%
\footline={\tenrm \ifodd\pageno 
            \hphantom{\folio}\hfil
            \lower12pt\hbox{\pdfrefximage\ZfSMEi}\unskip\hfil\folio \else
            \folio\hfil\lower12pt\hbox{\pdfrefximage\ZfSMEi}\unskip\hfil
            \hphantom{\folio}\fi}}
\fi\fi\fi\fi

% Formáty pre zarovnanie vľavo \fL, vpravo fR, na stred fC
% a do bloku s posledným riadkom vycentrovaným fX a \fXi (s odsadením) 
% viď https://petr.olsak.net/opmac-tricks.html#tabpmod
\def\fL{\rightskip=0pt plus 1fill}
\def\fR{\leftskip=0pt plus 1fill}
\def\fC{\leftskip=0pt plus 1fill \rightskip=0pt plus 1fill}
\def\fX{\leftskip=0pt plus1fil %\iindent plus1fil
        \rightskip=0pt plus-1fil %\iindent plus-1fil
        \parfillskip=0pt plus2fil}
\def\fXi{\leftskip=\iindent plus1fil
         \rightskip=\iindent plus-1fil
         \parfillskip=0pt plus2fil}

% Tlač pod/sekcie doplnená o zarovnanie textu vľavo (\fL)
\ifnum \SecLeft = 1
\def\printsec#1{\par \norempenalty-400 \bigskip
  {\fL \secfont \noindent \dotocnum
  {\thetocnum\quad}#1\nbpar}%
  \insertmark{#1}%
  \nobreak \remskip\medskipamount \firstnoindent
} 
\def\printsecc#1{\par \norempenalty-200 \medskip
  {\fL \seccfont \noindent \dotocnum
  {\thetocnum\quad}#1\nbpar}%
  \nobreak \remskip\medskipamount \firstnoindent
}
\fi

% Vytvorenie úrovne vyznačovania pre pod-pod-sekciu, text zarovnaný vľavo
% (bez číslovania a bez vytvárania záznamu pre obsah)
% viď https://petr.olsak.net/opmac-tricks.html#seccc
\def\seccc#1\par{\norempenalty-100 \medskip \noindent
{\ifnum \SecLeft = 1 \fL \fi \bf #1\nbpar}%
\nobreak \medskip \firstnoindent}

% Zmena nastavenia zarovnania riadkov textu podľa šírky strany
\ifdim\hsize < \MinHsize 
% bežný text bude zarovnaný vľavo, pokiaľ nie je nastavená väčšia rozťažnosť
\ifdim\emergencystretch>0pt \else  \rightskip = 0pt plus3em \fi
\else\ifdim\hsize < \MinHSize
\ifdim\emergencystretch>0pt \else
\emergencystretch=2em % bežný text zarovnaný do bloku s väčšou rozťažnosťou
\fi\fi\fi

% Formát pre autora na titulnej strane
\def\authorfont{\typobase\typoscale[\magstep4/\magstep4]\it}
\def\AuthorF#1{\nobreak\bigskip
{\fC\authorfont \noindent #1\par}%
}

% Formát pre zoznam digitalizátorov
\newdimen\rozmerD % rozmer stĺpca zoznamu digitalizátorov
\def\Digitalizator#1#2{\setbox0=\hbox{#1}% pre šírku textu 1. parametra
\rozmerD=\hsize
\advance \rozmerD by -\wd0 % zmenšené o šírku textu 1. stĺpca
\advance \rozmerD by -1em % zmenšené o cca medzeru medzi stĺpcami
\table{lp{\rozmerD\fL}}
{\kern-.5em#1 & #2 \cr}\bigskip} % \kern-.5em posun textu 1. stĺpca k okraju
                                 % tabuľky (potlačenie medzery) 

% Formát pre Copyright
\def\CopyR#1{\noindent{\fC #1\nbpar}}

% Formát pre Creative Commons License vycentruje poslený riadok, viď \fX
\def\CCL#1#2{\medskip\noindent #1 \par\nobreak
\noindent{\fX #2\par}
\medskip\noindent}

\def\ZFurl{% Vytlačí url autorských práv pre Zlatý fond SME
\noindent{\fC
\url{http://zlatyfond.sme.sk/dokument/autorske-prava/}\par}}

\def\TypoInfo{% Vytlačí info o TeX a použitom fonte
\ifdim\vsize > 110mm \vskip 12mm\noindent \else \medskip\noindent \fi
Sadzba programom \ulink[https://ctan.org]{\TeX} písmom \FontType .}

% Vytvorenie formátu pre jazykovú úpravu, text zarovnaný vľavo
\def\JazykUprava#1\par{\norempenalty-100 \bigskip \noindent
{\ifnum \SecLeft = 1 \fL \fi
\it #1\nbpar}\nobreak \bigskip \firstnoindent} 

% Formát pre Bibliogafické poznámky
\def\BibPoznamky#1\par{\norempenalty-100 \bigskip \noindent
{\fC\bf #1\nbpar}%
\nobreak \medskip \firstnoindent}

% Formát pre perex za názvom kapitoly vycentruje poslený riadok, viď TBN s. 234
% v prípade užšej strany ako \MinHsize bude perex zarovnaný vľavo
% po odkomentovaní %\catcode`\~=10 už nebude ~ nedeliteľná  medzera
\ifdim\hsize < \MinHsize %\catcode`\~=10 
\long\def\Perex#1#2{\noindent{\it #1\par}\medskip\noindent
{\sl #2}\medskip} \else
\long\def\Perex#1#2{\noindent{\fX\it #1\par}\medskip\noindent
{\sl #2}\medskip}\fi

% Formát pre motto za názvom kapitoly 
% v prípade užšej strany ako \MinHsize bude motto zarovnané vľavo
\ifdim\hsize < \MinHsize 
\long\def\Motto#1#2{{\noindent\fL\sl \MottoH \par
\smallskip\noindent #1\par}\medskip\noindent{\sl #2}\medskip} 
\else
\long\def\Motto#1#2{{\noindent\fL\sl \MottoH \par}
\Vers #1\par{\noindent\fL\sl #2\par}\medskip}
\fi

% Formát pre tabuľku Miery a váhy v závislosti od šírky strany
\def\TabMaV#1{%
\ifdim\hsize < \MinHsize 
\table{llp{42mm\fL}}{#1} %strana šírky 78 mm
\else\ifdim\hsize < \MinHSize 
\table{llp{52mm\fL}}{#1} %strana šírky 96 mm
\else
\table{lll}{#1} % strana väčšej šírky 
\fi\fi }

% Formát pre položku zoznamu \item obklopenú medzerami
\def\Itemm#1\par{\smallskip\item #1\smallskip}

% Oddeľovač hviezdička medzi odstavcami (2/2025 zmena smallskip->medskip):
\def\Hviezdicka{\nobreak\medskip\centerline{*}\medskip}

% Slovo Koniec a jeho formát pre celý riadok na konci kapitol
\def\Koniec#1\par{\nobreak\medskip\noindent{\it Koniec\ #1}}

% Formát pre verše
\def\nnarrower{\narrower\narrower}
\def\Vers#1\par{\ifdim\hsize < \MinHsize % pre formát ebsr bez odsadenia
{\removelastskip\smallskip\noindent
                 \it #1\smallskip}\else
{\removelastskip\smallskip\noindent\nnarrower
                 \it #1\smallskip}\fi
                }

% Definícia znaku stupeň s medzerou pre použitie mimo matematického módu 
\def\Deg{$\,^{\circ}\kern-.1em$}
% Definícia znaku stupeň bez medzery pre použitie mimo matematického módu 
\def\DEG{$^{\circ}$}

% Definícia zlomku pre použitie mimo matematického módu (\FR 1/2 je 1/2)
\def\FR#1/#2 {\ifnum\FRformat=0
$\,^{#1}\!/\!_{#2}\,$\else
$\,{#1\over#2}\,$\fi}
% Prepínač formátu zlomkov
\def\FRswitch{\ifnum\FRformat=0
\def\FRformat{1}\else\def\FRformat{0}\fi}

% Doplnenie znaku - na začiatok riadku v prípade delenia slov s týmto znakom
% TBN str. 217 (Rešení 3.)
\def\={\discretionary{-}{-}{-}}

% Potlačenie číslovania kapitol, sekcií a podsekcií
% https://petr.olsak.net/opmac-tricks.html#globnonum
\def\NumOff{\nonum \def\nonumfalse{\pagedepth=\pagedepth}}
\ifnum\NoNum = 1
\NumOff
% Namiesto názvu Kapitola štvotrec na pravej strane (pre nečíslované kapitoly)
\sdef{mt:chap:sk}{\hfill\vrule height12pt depth2pt width14pt}
\sdef{mt:chap:cs}{\hfill\vrule height12pt depth2pt width14pt}
\sdef{mt:chap:en}{\hfill\vrule height12pt depth2pt width14pt}
\fi
% Pre globálne obnovenie číslovania sa používa nasledovný odkomentovaný riadok:
%\newif\ifnonum \nonumfalse \def\nonum{\global\nonumtrue} \def\NoNum{0}

% Predefinovanie RS \chap \sec a \secc z OPmac pre písmenné označenie kapitol, 
% sekcí a podsekcí (napríklad A, B, C. ...) pomocou nových RS \chapnumf,
% \secnumf a \seccnumf.
%
% Použitie: pred označenie kapitoly/sekcie/podsekcie,  kde potrebujeme zmeniť
% jej formát  alebo číslo,  napísať  \def\chapnumf{p}/secnumf{p}/seccnumf{p},
% kde p je jedno veľké alebo malé písmeno bez diakritiky ([A-Z] alebo [a-z]),
% alebo  prirodzené číslo  (zložené len zo znakov [0-9]);  pri zápornom celom
% čísle sa budú používať  rímske číslice,  ktoré vyjadrujú  absolútnu hodnotu 
% záporného čísla.  Pre  doplnenie  bodky  za číslom  ako symbolu  pre radovú 
% číslovku  (často sa  používa  práve pri použití  rímskych číslic)  je možné 
% nastaviť pomocou \def\ChapRow{1}/SecRow{1}/SeccRow{1}  pre kapitoly/sekcie/
% /podsekcie a vypnúť pomocou ...{0}.
%
\def\Incrmc{1} \def\Decrmc{-1} 
\def\Incrms{1} \def\Decrms{-1}
\def\Incrmss{1} \def\Decrmss{-1}
\def\IfR#1{\ifnum#1<0% TBN s. 43, 365
\expandafter\uppercase\expandafter{\romannumeral-#1}\else\the#1\fi}
\def\IfRo#1{\ifnum#1<0% \IfRo sa používa pokiaľ namiesto \the \othe 
\expandafter\uppercase\expandafter{\romannumeral-#1}{}.\else\othe#1.\fi}
\newcount\charchap \newcount\charsec \newcount\charsecc
\def\chapnumf#1{%
    \ifcat A#1 \def\Incrmc{1} \def\Decrmc{-1}
               \chapnum=`#1 \advance\chapnum by\Decrmc \charchap=1 \else
    \ifnum#1 < 0 \def\Incrmc{-1} \def\Decrmc{1} \else
                 \def\Incrmc{1} \def\Decrmc{-1} \fi
    \chapnum=#1 \advance\chapnum by\Decrmc \charchap=0 \fi
}
\def\secnumf#1{%
    \ifcat A#1 \def\Incrms{1} \def\Decrms{-1}
               \secnum=`#1 \advance\secnum by\Decrms \charsec=1 \else
    \ifnum#1 < 0 \def\Incrms{-1} \def\Decrms{1} \else
                 \def\Incrms{1} \def\Decrms{-1} \fi
    \secnum=#1 \advance\secnum by\Decrms \charsec=0 \fi
}
\def\seccnumf#1{%
    \ifcat A#1 \def\Incrmss{1} \def\Decrmss{-1}
               \seccnum=`#1 \advance\seccnum by\Decrmss \charsecc=1 \else
    \ifnum#1 < 0 \def\Incrmss{-1} \def\Decrmss{1} \else
                 \def\Incrmss{1} \def\Decrmss{-1} \fi
    \seccnum=#1 \advance\seccnum by\Decrmss \charsecc=0 \fi
}
% + doplneneie definíce \ChapTitle do makra \chap
\eoldef\chap#1{\ifnonum\else \global\advance\chapnum by\Incrmc \fi
  \chaphook {\globaldefs=1 \secnum=0 \seccnum=0 \tnum=0 \fnum=0 \dnum=0%
             \charsec=0 \charsecc=0 \def\Incrms{1} \def\Decrms{-1}%
             \def\Incrmss{1} \def\Decrmss{-1}}\relax
  \ifnum\ChapRow=0 \def\ChRD{} \else \def\ChRD{{}.} \fi
  \ifnum\charchap=1%
  \edef\thechapnum{\char\the\chapnum}\let\thetocnum=\thechapnum\else
  \edef\thechapnum{\IfR{\chapnum}\ChRD}\let\thetocnum=\thechapnum\fi
  \edef\thesecnum{\IfRo{\chapnum}\IfR{\secnum}}%
  \def\dotocnumafter{\wtotoc0\bfshape{#1}}%
  \printchap{#1}\resetnonumnotoc
  \def\ChapTitle{#1} % doplnené do pôvodnej definície \chap
}
% Doplnenie začiatočnej nepárnej strany kapitoly  do makra \printchap
\def\printchap#1{\null\vfill\supereject
    \NewPageIfEven
  {\chapfont \noindent \mtext{chap} \dotocnum{\thetocnum}\par
   \nobreak\smallskip\noindent #1\nbpar}\mark{}%
   \nobreak\remskip\bigskipamount \firstnoindent
}
\eoldef\sec#1{\ifnonum\else \global\advance\secnum by\Incrms \fi
   \sechook {\globaldefs=1 \seccnum=0 \tnum=0 \fnum=0 \dnum=0%
             \charsecc=0 \def\Incrmss{1} \def\Decrmss{-1}}\relax
   \ifnum\SecRow=0 \def\SRD{} \else \def\SRD{{}.} \fi
   \ifnum\charchap=1  \ifnum\charsec=1% 
   \edef\thesecnum{\char\othe\chapnum.\char\the\secnum}%
   \let\thetocnum=\thesecnum\else
   \edef\thesecnum{\char\othe\chapnum.\IfR{\secnum}\SRD}
   \let\thetocnum=\thesecnum\fi
   \else
   \ifnum\charsec=1%
   \edef\thesecnum{\IfRo{\chapnum}\char\the\secnum}
   \let\thetocnum=\thesecnum\else
   \edef\thesecnum{\IfRo{\chapnum}\IfR{\secnum}\SRD}%
   \let\thetocnum=\thesecnum\fi\fi
   \def\dotocnumafter{\wtotoc1\rm{#1}}%   
   \printsec{#1}\resetnonumnotoc 
}
\eoldef\secc#1{\ifnonum\else \global\advance\seccnum by\Incrmss \fi
   \secchook {}\relax
   \ifnum\SeccRow=0 \def\SsRD{} \else \def\SsRD{{}.} \fi  
   \ifnum\charchap=1%
   \ifnum\charsec=1 \ifnum\charsecc=1% TBN, s. 47 (podmienka typu A and B)
   \edef\theseccnum{\char\othe\chapnum.\char\the\secnum.\char\the\seccnum}\else
   \edef\theseccnum{\char\othe\chapnum.\char\the\secnum.\IfR{\seccnum}\SsRD}%
   \fi\else
   \ifnum\charsecc=1%
   \edef\theseccnum{\char\othe\chapnum.\IfR{\secnum}.\char\the\seccnum}\else
   \edef\theseccnum{\char\othe\chapnum.\IfR{\secnum}.\IfR{\seccnum}\SsRD}%
   \fi\fi\else
   \ifnum\charsec=1 \ifnum\charsecc=1% 
   \edef\theseccnum{\IfRo{\chapnum}\char\the\secnum.\char\the\seccnum}\else
   \edef\theseccnum{\IfRo{\chapnum}\char\the\secnum.\IfR{\seccnum}\SsRD}%
   \fi\else
   \ifnum\charsecc=1%
   \edef\theseccnum{\IfRo{\chapnum}\IfRo{\secnum}\char\the\seccnum}\else
   \edef\theseccnum{\IfRo{\chapnum}\IfRo{\secnum}\IfR{\seccnum}\SsRD}\fi\fi\fi
   \let\thetocnum=\theseccnum
   \def\dotocnumafter{\wtotoc2\rm{#1}}%   
   \printsecc{#1}\resetnonumnotoc 
}

% Pomocné makrá s dvomi parametrami, ktoré sa použijú podľa podmienky: 
\def\IfEbook#1#2{%ak je výstupom e-book, použije sa 1.parameter, ak nie 2.
\ifnum\LogoZF=-1 #1 \else #2 \fi}
\def\IfHyperlinks#1#2{%ak aktívne hyperlinky, bude 1.parameter, ak nie 2.
\ifx\link\linkactive #1 \else #2 \fi}

% Pomocné makrá pre nastavenie záhlavia strany (viď TBN, str. 259-260)
\def\HeadSize{\thefontsize[10]}
\let\ChapTitle=\null
\def\HeadRule{\leaders\hrule height3pt depth-2.6pt\hfil} 
\def\MySpace{0pt} % východzia veľkost medzery medzi textom záhlavia a \HeadRule
\def\PraveZahlavie{\ifnum\NoNum=1 \else \def\MySpace{3pt}\fi
                   \HeadRule{\it\kern\MySpace\HeadSize\firstmark
                   \vphantom{\^Č}}}% 6.2.2025 doplnené \vphantom{\^Č}
\def\LaveZahlavie{\ifx\ChapTitle\null \else \def\MySpace{4pt}\fi
                 {\ifx\ChapTitle\firstmark \it \else \sl \fi
                  \HeadSize\vphantom{\^Č}% 6.2.2025 doplnené \vphantom{\^Č}
                  \ChapTitle\kern\MySpace \HeadRule}}
% Makro pre tlač záhlavia \TlacZahlavia na základe parametra \ZahlavieTlac
\def\TlacZahlavia{\ifnum\ZahlavieTlac = 1
\headline={\ifodd\pageno \PraveZahlavie \else \LaveZahlavie \fi}\fi}

% Nastavenie vkladania obrázkov rôzneho umiestnenia a rôznych formátov 
\def\InsPicc#1#2{\openin\testin=\ToWp #1 
\ifeof\testin\opwarning{The ‘#1’ file doesn’t exist..}\else
\closein\testin
\ifx^#2^\tmpdim=.4pt\else\tmpdim=#2\fi % TPP str. 71
\vbox{\hrule height\tmpdim\hbox{\vrule width\tmpdim
\inspic\ToWp #1 \vrule width\tmpdim}\hrule height\tmpdim}\fi
}
\def\InsPiccR#1#2{\openin\testin=\ZFroot #1
\ifeof\testin\opwarning{The ‘#1’ file doesn’t exist..}\else
\closein\testin
\ifx^#2^\tmpdim=.4pt\else\tmpdim=#2\fi % TPP str. 71
\vbox{\hrule height\tmpdim\hbox{\vrule width\tmpdim
\inspic\ZFroot #1 \vrule width\tmpdim}\hrule height\tmpdim}\fi
}
\def\InsPice#1#2{\openin\testin=\ToWp #1
\ifeof\testin\opwarning{The ‘#1’ file doesn’t exist..}\else
\closein\testin
\ifx^#2^\tmpdim=.4pt\else\tmpdim=#2\fi % TPP str. 71
\vbox{\hrule height\tmpdim\hbox{\vrule width\tmpdim
\epsfbox{\ToWp #1}\vrule width\tmpdim}\hrule height\tmpdim}\fi
}
\def\InsPiceR#1#2{\openin\testin=\ZFroot #1
\ifeof\testin\opwarning{The ‘#1’ file doesn’t exist..}\else
\closein\testin
\ifx^#2^\tmpdim=.4pt\else\tmpdim=#2\fi % TPP str. 71
\vbox{\hrule height\tmpdim\hbox{\vrule width\tmpdim
\epsfbox{\ZFroot #1}\vrule width\tmpdim}\hrule height\tmpdim}\fi
}
% Použitie rôznych verzií formátu obrázkov na základe parametra \ObrPdf
\def\FigF#1#2#3#4#5#6{%% vloží obrázok zarovnaný na ľavý okraj bez medzier
\ifnum\ObrPdf = 2      % vloží eps obrázok (pre dvi výstup)
\leftline{#6\InsPice{#5}{0pt}}
\else
\ifnum\ObrPdf = 1      % vloží čiernobiely obrázok (png, jpg, pdf)
\leftline{#4\InsPicc{#3}{0pt}}
\else                  % vloží farebný obrázok (png, jpg, pdf)
\leftline{#2\InsPicc{#1}{0pt}}%
\fi\fi
}
\def\FigFM#1#2#3#4#5#6{%% vloží obrázok zarovnaný na ľavý okraj s medzerami
\medskip
\ifnum\ObrPdf = 2       % vloží eps obrázok (pre dvi výstup)
\leftline{#6\InsPice{#5}{0pt}}
\else
\ifnum\ObrPdf = 1       % vloží čiernobiely obrázok (png, jpg, pdf)
\leftline{#4\InsPicc{#3}{0pt}}
\else                   % vloží farebný obrázok (png, jpg, pdf)
\leftline{#2\InsPicc{#1}{0pt}}
\fi\fi\medskip
}
\def\FigFC#1#2#3#4#5#6{%% vloží obrázok zarovnaný na stred šírky strany
\ifnum\ObrPdf = 2       % vloží eps obrázok (pre dvi výstup)
\centerline{#6\InsPice{#5}{0pt}}
\else
\ifnum\ObrPdf = 1       % vloží čiernobiely obrázok (png, jpg, pdf)
\centerline{#4\InsPicc{#3}{0pt}}
\else                   % vloží farebný obrázok (png, jpg, pdf)
\centerline{#2\InsPicc{#1}{0pt}}
\fi\fi
}
\def\FigFCframe#1#2#3#4#5#6#7#8#9{%% vloží obrázok s orámovaním na stred strany
\ifnum\ObrPdf = 2                  % vloží eps obrázok (pre dvi výstup)
\centerline{#8\InsPice{#7}{#9}}
\else
\ifnum\ObrPdf = 1                  % vloží čiernobiely obrázok (png, jpg, pdf)
\centerline{#5\InsPicc{#4}{#6}}
\else                              % vloží farebný obrázok (png, jpg, pdf)
\centerline{#2\InsPicc{#1}{#3}}
\fi\fi
}
\def\FigFCroot#1#2#3#4#5#6{%% vloží obrázok z ZFroot zarovnaný na stred strany
\ifnum\ObrPdf = 2           % vloží eps obrázok (pre dvi výstup)
\centerline{#6\InsPiceR{#5}{0pt}}
\else
\ifnum\ObrPdf = 1           % vloží čiernobiely obrázok (png, jpg, pdf)
\centerline{#4\InsPiccR{#3}{0pt}}
\else                       % vloží farebný obrázok (png, jpg, pdf)
\centerline{#2\InsPiccR{#1}{0pt}}
\fi\fi
}

% Generovanie okrajov strán (hoffset, voffset) pre patitul a titul, musia
% byť nazávislé od okrajov pre ostatné strany
\ifdim\pgwidth=0pt \pgwidth=210mm \pgheight=297mm \fi % nastaví sa na A4
\newdimen\HoffsetTitul \HoffsetTitul=\pgwidth
\advance\HoffsetTitul by -\hsize
\divide\HoffsetTitul  by 2 % rovnaký pravý a ľavý okraj (ako ich priemer)
\advance\HoffsetTitul by -1in
\newdimen\VoffsetTitul \newdimen\VsizeTitul \VsizeTitul=\vsize
\ifdim\pgwidth>176mm \VoffsetTitul=1in  \VsizeTitul=246mm \else % pre A4
\ifdim\pgwidth>148mm \VoffsetTitul=24mm \VsizeTitul=202mm \else % pre B5
\ifdim\pgwidth>105mm \VoffsetTitul=18mm \VsizeTitul=174mm \else % pre A5
\ifdim\pgwidth>78mm  \VoffsetTitul=3mm                    \else % pre eb, ebr
\ifdim\pgwidth=78mm  \VoffsetTitul=1mm                          % pre ebsr
\fi\fi\fi\fi\fi
\advance\VoffsetTitul by -1in

\def\PatitulStrana{% Vytvorenie strany pre patitul (1. generovaná strana)
\nopagenumbers % patitul až po obsah bez číslovaných strán
{\hoffset=\HoffsetTitul \voffset=\VoffsetTitul \vsize=\VsizeTitul
\ifnum\pageno>0 \pageno=-1 \fi % od patitulu až po obsah budú záporné strany
\KindleHR
\null\vskip .1\vsize % vertikálny posun o 0,1 výšky strany
\noindent
{\fC{\rm \Autor: {\it \Titul}}\par}
\bigskip
\FigFCroot%
{\LogoZfSMEbc}{}% farebná verzia (png)
{\LogoZfSMEbg}{}% č/b verzia (png)
{\LogoZfSMEbe}{}% verzia pre DVI=>PS (eps)
\ifnum\pageno=-1 \NewPage\NewPageIfEven \else \NewPage\NewPageIfOdd \fi
}}

\def\TitSkipMinus{% ak prvý (expandovaný) token \TitSkip nie je . potom je -
\if .\TitSkip \else -\fi}

\def\TitulMedzera{%
\ifdim\vsize > 100mm       % v prípade strany vyššej ako 100 mm,
\null\vskip \TitSkip\vsize % vertikálny posun o \TitSkip * výška strany 
\else \if -\TitSkipMinus   % alebo v prípade zápornej hodnoty \TitSkip,
\null\vskip \TitSkip\vsize % vertikálny posun o \TitSkip * výška strany
\fi\fi} 		   % (na cca zlatý rez)  

\def\TitulStrana{% Vytvorenie strany pre titul (2. generovaná strana)
\nopagenumbers % titulá strana až po obsah bez číslovaných strán
\ifnum\pageno>0 \pageno=-1 \fi % od titulu až po obsah budú záporné strany 
{\hoffset=\HoffsetTitul \voffset=\VoffsetTitul \vsize=\VsizeTitul
\KindleHR
\TitulMedzera    		% posun titulu na cca zlatý rez
\csname\string\tit:M\endcsname{\Titul} % opmac-d.pdf, str. 16
\AuthorF \Autor\par
\ifnum\pageno=-1 \NewPage\NewPageIfEven \else \NewPage\NewPageIfOdd \fi
}}

% Závesná interpunkcia -- viď TPP str. 110
\newcount\tmpnum \newdimen\tmpdim % deklarácie z OPmac
\def\setcharcode#1#2#3{% #1: \whatcode, #2: char, #3: factor
    \setcharcodeforfont \tenrm #1#2{#3}%
    \setcharcodeforfont \tenbf #1#2{#3}%
    \setcharcodeforfont \tenit #1#2{#3}%
    \setcharcodeforfont \tenbi #1#2{#3}%
}
\def\setcharcodeforfont#1#2#3#4{%
    \setbox0=\hbox{#1#3}\tmpdim=#4\wd0 \tmpnum=\tmpdim
    \tmpdim=\fontdimen6#1% velikost 1em
    \divide\tmpdim by50 \multiply\tmpnum by20 \divide\tmpnum by\tmpdim
    #2#1\ifcat#3\clqq\else\expandafter`\fi#3=\tmpnum
}
\def\ZavesnaIterpunkcia{
    \chardef\hyph=\hyphenchar\font
    \setcharcode \lpcode\clqq {.9}
    \setcharcode \rpcode\crqq {.9}
    \setcharcode \rpcode .    {1}
    \setcharcode \rpcode ,    {1}
    \setcharcode \lpcode ,    {1}  %doplnené , pre funkciu ,úvodzovky‘
    \setcharcode \rpcode\elq  {1}  %doplnené ‘ pre funkciu ,úvodzovky‘
    \setcharcode \lpcode\elqq {.9} %doplnená “americká” ľavá úvodzovka
    \setcharcode \rpcode\erqq {.9} %doplnená “americká” pravá úvodzovka
    \setcharcode \lpcode\elq  {1}  %doplnená ‘americká’ ľavá úvodzovka
    \setcharcode \rpcode\erq  {1}  %doplnená ‘americká’ pravá úvodzovka
    \setcharcode \rpcode\hyph {.9} %zmenené z pôvodnej hodnoty 0.7 na 0.9 
    \setcharcode \lpcode\hyph {.9} %doplnené o \lpcode\hyph pre použitie v \=
    %% pre "veľké" interpunkčné znamienka sa závesná interpunkcia nepoužíva,
    %% ale po odkomentovaní je možné použiť aj tieto znaky:
    %\setcharcode \lpcode\flqq {.9} %FR <<úvodzovky, doplnené voči TPP
    %\setcharcode \rpcode\flqq {.9} %FR >>úvodzovky, doplnené voči TPP
    %\setcharcode \lpcode\frqq {.9} %FR úvodzovky>>, doplnené voči TPP
    %\setcharcode \rpcode\frqq {.9} %FR úvodzovky<<, doplnené voči TPP
    %\setcharcode \rpcode ;    {1}  %doplnené voči TPP
    %\setcharcode \rpcode :    {1}  %doplnené voči TPP
    %\setcharcode \rpcode !    {1}  %doplnené voči TPP
    %\setcharcode \rpcode ?    {.8} %doplnené voči TPP
    %% koniec "veľkých" interpunkčných znamienok
    \pdfprotrudechars=2 
}
\ifdim\pgwidth=0pt \pgwidth=210mm \fi % \pgwidth sa nastaví na hodnotu pre A4
\newdimen\OkrajeH \OkrajeH=\pgwidth \advance\OkrajeH by -\hsize 
\ifdim\OkrajeH > \MinOkraj % pre úzke okraje sa závesná inerpunkcia nepoužije
    \ZavesnaIterpunkcia
\fi
%% Definovanie vypnutia a zapnutia ľavých znakov závesnej interpuncie pre
%% odstavec s \Capinsert, viď súbor ../Reuss_Gustav/G_EAAAA.tex, kap. Afrika
\ifdim\OkrajeH > \MinOkraj %vypnutie ľavých znakov (pre odstavec s \Capinsert)
\def\HPleftOff{\setcharcode \lpcode\clqq {0}
           \setcharcode \lpcode ,    {0}
           \setcharcode \lpcode\elqq {0}
           \setcharcode \lpcode\elq  {0}
           \setcharcode \lpcode\hyph {0}
           %\setcharcode \lpcode\flqq {0} %pripravené pre FR ľavú úvodzovku
           %\setcharcode \lpcode\frqq {0} %pripr. pre FR obrátenú úvodzovku
           }
\else \def\HPleftOff{} \fi 
\ifdim\OkrajeH > \MinOkraj %zapnutie ľavých znakov (po odstavci s \Capinsert)
\def\HPleftOn{\setcharcode \lpcode\clqq {.9}
          \setcharcode \lpcode ,    {1}
          \setcharcode \lpcode\elqq {.9}
          \setcharcode \lpcode\elq  {1}
          \setcharcode \lpcode\hyph {.9}
          %\setcharcode \lpcode\flqq {.9} %pripravené pre FR ľavú úvodzovku 
          %\setcharcode \lpcode\frqq {.9} %pripr. pre FR obrátenú úvodzovku
          }
\else \def\HPleftOn{} \fi

% Rozdelenie inexov do oddielov s nadpisom podľa písmena
% https://petr.olsak.net/opmac-tricks.html#currchar
% nevhodné pre RS, pred položku s RS vložiť napr. \null
% vo vim :%s/{{\\tt/{\\null{\\tt/gc
\ifnum\IndexChar = 1
\def\lastchar{\null} % doplnené voči opmac-tricks.html#currchar 
\def\everyii{\expandafter\makecurrchar\currii\end
    \ifx\currchar\lastchar \else
    \if\currchar\lastchar \else % doplnené voči opmac-tricks.html#currchar
       \bigskip{\typosize[12/16]\bf\currchar}\par\nobreak\medskip\noindent
       \let\lastchar=\currchar
    \fi\fi 
} 
\def\makecurrchar#1#2\end{\uppercase{\def\currchar{#1}}}
\fi
% Klikacie odkazy na čísla strán v indexových položkách
% https://petr.olsak.net/opmac-tricks.html#klikpgr
% určené len pre malý počet indexových položiek
\ifnum\IndexLink = 1
\def\printiipages#1&{\usepglinks#1\relax\par}
\def\usepglinks{\afterassignment\usepglinksA \tmpnum=}
\def\usepglinksA{\pglink{\the\tmpnum}\futurelet\next\usepglinksB}
\def\usepglinksB{\ifx\next\relax \def\next{}\else
    \ifx\next,\def\next,{, \afterassignment\usepglinksA \tmpnum=}\else
    \ifx\next-\def\next--{--\afterassignment\usepglinksA\tmpnum=}%
    \fi\fi\fi\next}
\fi

%% https://petr.olsak.net/opmac-tricks.html#capinsert 
%% nastavenie jednotiek (\ptex a \VshiftFont sa nastaví už pri konfigurácii)
\newdimen\ptem       \ptem=.1em          % jednotka závislá na em 
                     \ptex=.1ex          % jednotka závislá na ex 
%\newdimen\ptex      \ptex=.1ex          % jednotka závislá na ex 
\newdimen\Capsize    \Capsize=70\ptex    % požadovaná veľkosť (2 riadky) 
\newdimen\Capabove   % horný presah cez účarie, hodnota je v \declCap
\ifnum\CapNumLines=3
\Capsize=116\ptex                        % požadovaná veľkosť (3 riadky)
\VshiftFont=1.7\VshiftFont               % pomer veľkostí 116/70 v \ptex
\fi
\newdimen\Capafter   \Capafter=1\ptem    % medzera za písmenom 
%\newdimen\VshiftFont \VshiftFont=0\ptex % vertikálny posun pre písmeno
%% vlastné makrá
\def\declCap #1#2{\sxdef{cap:=#1}{#2}}
\def\Capinsert{\def\leftCapmaterial{}\futurelet\next\CapinsertA} 
\def\CapinsertA{\ifx\next\bgroup \expandafter\CapinsertB \else 
\expandafter\CapinsertC \fi} 
\def\CapinsertB #1{\def\leftCapmaterial{#1}\CapinsertC} 
\def\CapinsertC #1{\par 
  \isdefined{cap:=#1}\iftrue \edef\tmp{\csname cap:=#1\endcsname}% 
                     \else   \edef\tmp{\csname cap:=default\endcsname}\fi 
  \setbox0=\hbox{{\thefontsize[\expandafter\ignorept\the\Capsize]\Capprefix#1}\kern\Capafter}% 
  \expandafter \CapinsertD \tmp,,% 
  \noindent\kern-\firstlineindent \rlap{\kern-\protrudeCap\ptem\llap{\leftCapmaterial}% 
  \advance\Capabove by \VshiftFont      \vbox to0pt{\kern-\Capabove\box0\vss}}% 
  \kern\firstlineindent 
} 
\def\CapinsertD #1:#2;{\tmpnum=1 \let\firstlineindent=\undefined 
   \Capabove=#2\ptex \def\parshapeparams{}\def\protrudeCap{#1}\CapinsertE} 
\def\CapinsertE #1,{\ifx,#1,\parshape =\tmpnum \parshapeparams 0pt \hsize 
  \else 
     \advance\tmpnum by1 
     \tmpdim=\wd0 \advance\tmpdim by-#1\ptem \advance\tmpdim by-\protrudeCap\ptem 
     \edef\parshapeparams{\parshapeparams\the\tmpdim}% 
     \ifx\firstlineindent\undefined \let\firstlineindent\parshapeparams \fi 
     \advance\tmpdim by-\hsize \tmpdim=-\tmpdim 
     \edef\parshapeparams{\parshapeparams\the\tmpdim}% 
     \expandafter \CapinsertE \fi 
} 
\def\hboxshift#1#2{\vbox to0pt{\vss\hbox{#2}\kern-#1}}
%% Deklarácia k jednotlivým písmenám v tvare: 
% \declCap písmeno {presah doľava:hore; korekcia prvého riadka, druhého, ...}
% (potrebné postupne doplniť aj o ďalšie písmená podľa potreby)
\declCap {default} {0:20;0,0}  % východzia hodnota pre nedeklarované písmená 
\declCap A {1:20;6,-2} 
\declCap B {2:20;1,-1}
\declCap D {1:20;0,0}
\declCap Ď {1:30;0,0}
\declCap E {1:20;0,-1}
\declCap I {1:20;0,-1}
\declCap J {1:21;2,1}
\declCap K {1:20;1,-2} 
\declCap L {0:20;9,0}
\declCap Ľ {0:20;1,-2}
\declCap N {2:20;0,-1}
\declCap P {1:20;0,6} 
\declCap U {1:21;2,1}
\declCap V {2:21;0,6}
\declCap W {3:20;0,6}
\declCap Z {1:20;1,-1}
\ifdim\hsize < \MinHsize \declCap Č {-1.2:31;0,-2} \else
\declCap Č {1:31;0,-2} \fi
\ifx\FontType\CMCS \declCap Ľ {0:30;2,-2} \else \ifdim\hsize < \MinHsize
\declCap Č {-1.2:36;0,-2} \else \declCap Č {1:36;0,-2}
\fi\fi
\ifnum\CapNumLines=3
\declCap {default} {0:24;0,0,0}  % východzia hodnota pre nedeklarované písmená 
\declCap A {1:24;6,2,-2}
\declCap B {2:24;2,-1,-1}
\declCap Č {1:42;0,-2,-2}
\declCap D {1:24;2,-1,3} 
\declCap Ď {1:40;2,-1,3}
\declCap E {1:24;0,-1,-1}
\declCap I {1:24;0,-1,-1}
\declCap J {1:25;2,0,0}
\declCap K {1:24;1,3,-3}
\declCap L {0:24;9,0,0}
\declCap Ľ {0:24;2,-2,-2}
\declCap N {1:24;2,0,0}
\declCap P {1:24;0,0,6}
\declCap U {1:25;2,1,1}
\declCap V {2:24;0,3,6}
\declCap W {3:24;0,4,6} 
\declCap Z {1:24;1,-1,-1}
\ifx\FontType\CMCS 
\declCap Ľ {0:41;2,-2,-2}
\else
\declCap Ď {1:46;2,-1,3}
\fi\fi
%% V prípade ak je 1. odstavec krátky a ďaľšie odstavce môžu zasiahnuť do 
%% zväčšeného písmena, je potrebné ich potrebné zrušiť, ale v prípade ak sa
%% zväčšené písmeno nepoužije, vykonať odriadkovanie a odsadenie.
\def\CapIskip{} \def\CapIIskip{} 
\def\HIleftOff{} \def\HIIleftOff{} \def\HIleftOn{} \def\HIIleftOn{}
\ifnum\IfCapInsert>0 \let\nli\nl \let\CapIskip=\smallskip 
\let\HIleftOff=\HPleftOff \let\HIleftOn=\HPleftOn \else 
\def\nli{\nl\indent}\fi
\ifnum\IfCapInsert>1 \let\nlii\nl \let\CapIIskip=\smallskip 
\let\HIIleftOff=\HPleftOff \let\HIIleftOn=\HPleftOn \else 
\def\nlii{\nl\indent}\fi
%% Zapnutie funkcie \Capinsert na základe hodnoty \IfCapInsert
\ifnum\IfCapInsert<1 \def\CapInsert{\noindent} \def\CapIInsert{\noindent} \else 
\ifnum\IfCapInsert>1 \let\CapInsert=\Capinsert \let\CapIInsert=\Capinsert  
\else                \let\CapInsert=\Capinsert \def\CapIInsert{\noindent}\fi\fi
%% Ak je potrebné umiestniť úvodzovku pred zväčšené písmeno vľavo dole 
%% pozn. \kern0pt bolo pôvodne \kern-2pt, ktoré sa presunulo do parametra
%% \CapIInsert pre súbor ../Reuss_Gustav/E_Hviezdy.tex
\def\LeftCapIMaterial#1{{\ifnum\IfCapInsert<1 {#1}\else
\thefontscale[\magstep2]\localcolor\Grey\hboxshift{34\ptex}{#1}\kern0pt\fi}}
\def\LeftCapIIMaterial#1{{\ifnum\IfCapInsert<2 {#1}\else
\thefontscale[\magstep2]\localcolor\Grey\hboxshift{34\ptex}{#1}\kern0pt\fi}}
\ifnum\CapNumLines=3
\def\LeftCapIMaterial#1{{\ifnum\IfCapInsert<1 {#1}\else
\thefontscale[\magstep3]\localcolor\Grey\hboxshift{62\ptex}{#1}\kern0pt\fi}}
\def\LeftCapIIMaterial#1{{\ifnum\IfCapInsert<2 {#1}\else
\thefontscale[\magstep3]\localcolor\Grey\hboxshift{62\ptex}{#1}\kern0pt\fi}}
\fi


\endinput

